\section{Detection Cost (O1)}
\label{sec:detection}

Proposition~\ref{prop:cost} predicts three things: the passive cost grows as $1/p^2$; the witness's cost is set by the epoch volume until $p$ approaches the floor $f$, where it has a wall of its own; and a zero-byte counter pair with the same information as the witness is separated from it only by read skew. This section measures all three. Packets to detect are counted from the fault's onset. In Figure~\ref{fig:scaling} and Table~\ref{tab:detection}, we show the packets each arm needed to flag the faulty spine and the fraction of seeds in which it did so, for eight loss rates from 1.5 percent down to the floor of $10^{-5}$.

\begin{figure}[t]
\centering
\includegraphics[width=\columnwidth]{scaling_curve.pdf}
\caption{Median post-onset packets to detect a single-spine grayhole against loss rate $p$, over 50 seeds, $k=8$ spines, and a budget of 80 post-onset epochs of 2 million packets (dotted ceiling); shaded bands are interquartile ranges. Filled markers are cells with an action rate of 1.00; hollow markers are annotated with the fraction of seeds that detected, and their median is over those seeds only; a cross on the ceiling marks a cell in which no seed detected. The ledger's cost is one epoch until the loss rate nears the floor, then rises to its own wall; the passive detectors' cost grows toward the ceiling and crosses it two decades earlier.}
\label{fig:scaling}
\end{figure}

\begin{table}[t]
\caption{Detection of a single-spine grayhole within a budget of 80 post-onset epochs (160 million packets), 50 seeds per cell: action rate (Act.) and, over the seeds that detected, the median post-onset packets to the first correct flag (Pkts). The 95 percent Wilson interval is [0.93, 1.00] for every 1.00, [0.65, 0.87] for 0.78, [0.26, 0.52] for 0.38, and [0.00, 0.07] for every 0.00. False positives were 0.00 for every arm at every rate.}
\label{tab:detection}
\centering
\footnotesize
\begin{tabular}{@{}lcccccc@{}}
\toprule
 & \multicolumn{2}{c}{Ledger} & \multicolumn{2}{c}{SprayCheck-Z} & \multicolumn{2}{c@{}}{FlowPulse-$\theta$} \\
\cmidrule(lr){2-3}\cmidrule(lr){4-5}\cmidrule(l){6-7}
$p$ & Act. & Pkts & Act. & Pkts & Act. & Pkts \\
\midrule
1.5\%     & 1.00 & 2\,M & 1.00 & 4\,M & 1.00 & 2\,M \\
1.0\%     & 1.00 & 2\,M & 1.00 & 6\,M & 1.00 & 2\,M \\
0.5\%     & 1.00 & 2\,M & 1.00 & 12\,M & 0.38 & 68\,M \\
$10^{-3}$ & 1.00 & 2\,M & 0.78 & 94\,M & 0.00 & -- \\
$10^{-4}$ & 1.00 & 2\,M & 0.00 & -- & 0.00 & -- \\
$5\!\times\!10^{-5}$ & 1.00 & 6\,M & 0.00 & -- & 0.00 & -- \\
$2\!\times\!10^{-5}$ & 1.00 & 36\,M & 0.00 & -- & 0.00 & -- \\
$10^{-5}$ ($=f$) & 0.00 & -- & 0.00 & -- & 0.00 & -- \\
\bottomrule
\end{tabular}
\end{table}

\textbf{The ledger detects on one epoch until the floor is near.} From 1.5 percent down to $10^{-4}$ the ledger flagged the faulty spine in every seed on the first epoch after onset, at a cost of 2 million fleet packets with an interquartile range of 2 to 2 million, widening to 2 to 4 million at $10^{-4}$. The reason is that one epoch carries 250 thousand packets per spine, which exceeds the $s^2 f/p^2$ of Proposition~\ref{prop:cost} at every rate in that range, so the first epoch's drops can already stand above the floor. Below $10^{-4}$ the wall appears where the proposition puts it. At $5\times10^{-5}$ the median cost rose to 6 million packets and at $2\times10^{-5}$ to 36 million, still with an action rate of 1.00 and no false positive. At the floor itself, $10^{-5}$, no seed detected, because a link that loses at the background rate is not a fault by the ledger's own definition. The flat portion of the curve is therefore a statement about the epoch volume relative to the floor, and the paper no longer calls it flat without that qualification.

\textbf{The passive cost grows as $1/p^2$ and the action rate collapses.} SprayCheck-Z needed 4, 6, and 12 million packets at 1.5, 1.0, and 0.5 percent, i.e., two, three, and six epochs. At $10^{-3}$ its median rose to 94 million packets, 47 times the ledger's, and that median is over the 78 percent of seeds that detected at all; its interquartile range widened from 4 to 4 million at 1.5 percent to 62 to 128 million at $10^{-3}$. At $10^{-4}$ and below no seed detected within the budget. FlowPulse-$\theta$ failed one order of magnitude earlier: it matched the ledger at 1.5 and 1.0 percent, detected in 38 percent of seeds at 0.5 percent with a median of 68 million packets, and detected nothing at $10^{-3}$ and below. Both passive arms had a false-positive rate of 0.00 at every rate, so the collapse is a loss of sensitivity, not a trade against false alarms. The separation is visible inside the budget: the 22 percent SprayCheck missed at $10^{-3}$ and its whole column at $10^{-4}$ are the same $s^2/p^2$ wall pushed past 160 million packets. A larger budget would move the passive floor to a rarer $p$ but cannot change the exponent; the ledger's curve would not move at all, because within its flat range it detects on the first epoch. This is a statement about cost, not capability: a passive detector with an unbounded budget can eventually detect at $10^{-4}$, and nothing here says otherwise.

\begin{figure}[t]
\centering
\includegraphics[width=\columnwidth]{fig_floor.pdf}
\caption{The ledger's wall tracks the floor: median epochs to detect against the loss rate as a multiple of the background floor $f$, for three floors, 50 seeds per cell. A cross at the top marks a cell in which no seed detected within 80 epochs; a hollow marker is annotated with its action rate. At every floor the ledger detects nothing at $p=f$, detects in every seed from about $5f$ upward, and reaches one epoch by $10f$ to $50f$.}
\label{fig:floor}
\end{figure}

\textbf{The wall moves with the floor.} Figure~\ref{fig:floor} repeats the ledger arm at three background floors, $10^{-6}$, $10^{-5}$, and $10^{-4}$, over $p$ from $f$ to $50f$. The curves align when plotted against $p/f$. At $p=f$ nothing is detected at any floor. At $2f$ the action rate is 1.00 at the two higher floors, in a median of 18 and 8.5 epochs after onset, and 0.18 at $10^{-6}$, where the floor produces half a drop per epoch and the separation is the sparse Poisson case of the proposition. From $5f$ upward every seed detects, in a median of 14, 3, and 2 epochs after onset at floors of $10^{-6}$, $10^{-5}$, and $10^{-4}$, falling to one epoch by $50f$, $10f$, and $10f$. Consequently, we expect the loss rate at which the witness stops paying to be not a fixed number but a small multiple of whatever the fabric's background floor is, which production floors can put anywhere from $10^{-6}$ to $10^{-4}$~\cite{corropt}. SprayCheck's action rate at a fixed $p$ barely changes with the floor, since its noise is spray noise rather than floor noise: 0.16 at $p=5\times10^{-4}$ over a $10^{-5}$ floor and 0.06 over $10^{-4}$.

\begin{figure*}[t]
\centering
\includegraphics[width=\textwidth]{fig_counterpair.pdf}
\caption{The zero-byte counter pair against controller read skew, 50 seeds per cell, for three loss rates. (a) Fraction of seeds in which a healthy spine was flagged. (b) Fraction of seeds in which the faulty spine was flagged. At $\sigma=0$ the arm is the ledger's information read at one instant and matches the ledger exactly; the vertical line is the best case measured on the switch, a 2.6 ms pairwise read against a 100 ms epoch.}
\label{fig:counterpair}
\end{figure*}

\textbf{What the two bytes buy: in-band alignment.} Figure~\ref{fig:counterpair} answers the question a per-link counter pair raises. At $\sigma=0$ the counter pair reproduced the ledger cell for cell, 2 million packets and no false positive at every rate, which confirms that Equation~\eqref{eq:loss} contains no information a counter pair lacks. The two arms separate as soon as the two registers are read at different instants. At $\sigma=10^{-4}$, ten microseconds against a 100 ms epoch, the counter pair still flagged the faulty spine in every seed. It also flagged a healthy spine in 28 percent of seeds at 0.5 percent loss and above, in 94 percent at $10^{-3}$, and in every seed at $10^{-4}$. At $\sigma=10^{-3}$ and above, including the measured best case of $\sigma=0.026$, it flagged a healthy spine in every seed at every rate. At the floor, where there is no fault, it still raised an alarm in 62 percent of seeds at $\sigma=10^{-4}$ and in every seed above that. The reason is that a read skew of $\sigma$ on a link carrying $\lambda$ packets per epoch adds phantom loss with standard deviation near $\lambda\sigma/\sqrt{6}$. That is about 10 packets at $\sigma=10^{-4}$ and 2\,700 at $\sigma=0.026$, against a true signal of 25 packets at $10^{-4}$ and a floor of 2.5. The arm does not lose the faulty link; it loses the ability to say which links are healthy, which we believe is the failure a repair action can least afford. Section~\ref{sec:localization} shows the consequence for localization. The witness reads both registers on one switch in one batched operation, so its skew is zero by construction. That, and not any inference, is what two bytes per packet buy. A counter pair would need to match it by reading its two registers within about ten microseconds of each other across two switches, four orders of magnitude tighter than the 350 ms census and two orders tighter than the 2.6 ms pairwise read that we measured.

\textbf{Two disclosures.} SprayCheck's sensitivity is calibrated under its paper's Poisson model at 2.5 million packets per spine, the cumulative count its Z-test reaches after ten epochs; its real join-the-shortest-queue spraying has less noise, and a sensitivity tuned to it would lower its curve by a constant factor without changing the exponent. The conditional medians of the hollow points in Figure~\ref{fig:scaling} are survivorship-biased, which is why the figure annotates each with its action rate and why we read the action-rate collapse, not the slope of a censored median, as the load-bearing signal.
